\documentclass[11pt]{article}
\usepackage{amsmath,amssymb,amsthm}
\usepackage{mathtools}
\usepackage{hyperref}
\usepackage[margin=1in]{geometry}
\usepackage[ruled,vlined,linesnumbered]{algorithm2e}
\usepackage{tikz}
\usetikzlibrary{arrows.meta,positioning}

\newtheorem{theorem}{Theorem}[section]
\newtheorem{lemma}[theorem]{Lemma}
\newtheorem{proposition}[theorem]{Proposition}
\newtheorem{corollary}[theorem]{Corollary}
\newtheorem{definition}[theorem]{Definition}
\newtheorem{remark}[theorem]{Remark}
\newtheorem{axiom_}{Axiom}[section]

\DeclareMathOperator*{\argmin}{arg\,min}
\DeclareMathOperator{\proj}{proj}
\DeclareMathOperator{\HM}{HM}
\DeclareMathOperator{\AM}{AM}
\DeclareMathOperator{\tr}{tr}
\DeclareMathOperator{\diag}{diag}

\title{Cooperative Policy Gradient with Lossy Communication\\
in General Stochastic Games}
\author{Eugene Shcherbinin}
\date{March 2026}

\begin{document}
\maketitle

\begin{abstract}
We extend the evidence-weighted policy gradient (EW-PG) and opponent-shaping
PG (LOLA-PG) frameworks to cooperative settings where agents form coalitions
and communicate policies through lossy channels. Three key contributions:
(1)~We formalize the \emph{self-knowledge problem}: an agent's ability to
communicate its policy is bounded by its evidence quality, connecting the
$\mathcal{O} \to \Pi$ inference gap to information-theoretic communication
limits. (2)~We show that communication introduces a bias decomposable into
self-knowledge loss and channel loss, with convergence preserved under
annealed cooperation. (3)~We prove that coalition formation enlarges the
basin of attraction through a cooperative reinforcement mechanism dual to
LOLA's spectral reinforcement. The evidence weight $w_i$ serves triple
duty: gradient quality, self-knowledge, and communication quality --- all
manifestations of the same underlying information content. This completes
the formalization of all five components of the $\Omega$-gradient.
\end{abstract}


\section{The Self-Knowledge Problem}

\subsection{The $\mathcal{O} \to \Pi$ Gap}

In the standard PG framework, agent~$i$ has policy $\pi_i \in \Pi$ and
computes gradient estimates $\hat{v}_i$ to improve it. The implicit
assumption is that the agent \emph{knows} its own policy. But in the
$\Omega$-framework, agents operate in observation space $\mathcal{O}$
and their policy is an implicit function of parameters~$\theta_i$:
\[
  \pi_i = \pi_{\theta_i} : \mathcal{S} \to \Delta(\mathcal{A}_i).
\]
The agent does not have direct introspective access to the function
$\pi_{\theta_i}$. It can only \emph{infer} a self-model $\hat{\pi}_i$
from its observations:
\[
  \hat{\pi}_i = f(o_1, a_1, r_1, \ldots, o_t) \approx \pi_i.
\]

\begin{definition}[Self-knowledge loss]
The self-knowledge loss of agent~$i$ is:
\begin{equation}\label{eq:self-knowledge}
  L_{\text{self}}^{(i)} \;=\; \mathbb{E}\bigl[\|\pi_i - \hat{\pi}_i\|^2\bigr].
\end{equation}
This measures the gap between ``what I do'' and ``what I think I do.''
\end{definition}

\begin{remark}[The ``vibing'' agent]
An agent that performs well through habitual or intuitive action --- what
we call ``vibing'' --- has a good policy~$\pi_i$ but a poor self-model
$\hat{\pi}_i$. In flow-state terms (Parvizi-Wayne et al.\ 2024), the
explicit self-model is attenuated while implicit performance is optimal.
In Polanyi's terms, the agent has \emph{tacit knowledge} that it cannot
articulate.
\end{remark}

\subsection{Evidence as Self-Knowledge}

The self-knowledge loss is related to the evidence variance $V_i$ from
the EW-PG framework. Both are driven by the same information deficit:
insufficient experience to resolve the $\mathcal{O} \to \Pi$ mapping.

\begin{axiom_}[Self-knowledge--evidence relationship]\label{ax:self-ev}
For each agent~$i$ with evidence variance $V_i > 0$:
\begin{equation}\label{eq:self-ev}
  L_{\text{self}}^{(i)} \;\leq\; V_i.
\end{equation}
\end{axiom_}

\noindent\textbf{Justification.}
By the data processing inequality,
$I(\pi_i;\, \hat{\pi}_i) \leq I(\pi_i;\, \mathcal{O}_i)$.
The evidence variance $V_i$ is inversely related to $I(\pi_i;\, \mathcal{O}_i)$:
high variance means the observation history is insufficient to pin down
the value function, and \emph{a fortiori} insufficient to pin down the policy.
The bound \eqref{eq:self-ev} makes this precise.


\section{The Communication Channel}

\subsection{End-to-End Communication}

When agent~$i$ communicates its policy to teammate~$j$, the information
passes through three stages:
\begin{equation}\label{eq:comm-chain}
  \underbrace{\pi_i}_{\text{true}} \;\xrightarrow{\text{self-model}}\;
  \underbrace{\hat{\pi}_i}_{\text{inferred}} \;\xrightarrow{\text{channel}}\;
  \underbrace{\tilde{\pi}_i^{(j)}}_{\text{received}}
\end{equation}

\begin{definition}[Communication loss]
The total communication loss for agent~$j$'s model of agent~$i$:
\begin{equation}
  L_{\text{comm}}^{(i \to j)} \;=\; \mathbb{E}\bigl[\|\pi_i - \tilde{\pi}_i^{(j)}\|^2\bigr].
\end{equation}
\end{definition}

\begin{theorem}[Loss decomposition]\label{thm:loss-decomp}
The communication loss decomposes as:
\begin{equation}\label{eq:loss-decomp}
  L_{\text{comm}}^{(i \to j)} \;\leq\; 2\,L_{\text{self}}^{(i)} + 2\,L_{\text{channel}}^{(i \to j)}
\end{equation}
where $L_{\text{channel}}^{(i \to j)} = \mathbb{E}[\|\hat{\pi}_i - \tilde{\pi}_i^{(j)}\|^2]$
is the pure channel distortion.
\end{theorem}

\begin{proof}
Write $\pi_i - \tilde{\pi}_i^{(j)} = (\pi_i - \hat{\pi}_i) + (\hat{\pi}_i - \tilde{\pi}_i^{(j)})$.
By $\|a + b\|^2 \leq 2\|a\|^2 + 2\|b\|^2$ (which follows from $(|a| - |b|)^2 \geq 0$),
we get \eqref{eq:loss-decomp}.
\end{proof}

\begin{corollary}[Self-knowledge bottleneck]\label{cor:bottleneck}
Even with a perfect channel ($L_{\text{channel}} = 0$):
\begin{equation}
  L_{\text{comm}}^{(i \to j)} \;\leq\; 2\,V_i.
\end{equation}
\emph{An agent that doesn't know its own policy cannot communicate it.}
\end{corollary}

\subsection{Rate-Distortion with Self-Knowledge Floor}

The communication problem connects to Shannon's rate-distortion theory.
Agent~$i$ encodes $\hat{\pi}_i$ at rate $R$ (bits per policy dimension).
The channel's rate-distortion function $D_{\text{ch}}(R)$ gives the minimum
achievable channel distortion at rate~$R$.

The total distortion has a floor:
\begin{equation}\label{eq:rate-distortion}
  D_{\text{total}}(R) \;\geq\; \max\bigl(D_{\text{ch}}(R),\; L_{\text{self}}\bigr).
\end{equation}

\begin{remark}[Bandwidth vs.\ evidence]
When $L_{\text{self}} > D_{\text{ch}}(R)$, the bottleneck is self-knowledge,
not bandwidth. Increasing channel capacity does not help --- the agent needs
more \emph{evidence}, not more bandwidth. The crossover rate
$R^* : D_{\text{ch}}(R^*) = L_{\text{self}}$ defines sufficient bandwidth
for the agent's evidence level.
\end{remark}


\section{Coalition Formation}

\subsection{Coalition Rationality}

\begin{definition}[Coalition value]
In an $N$-player stochastic game, the coalition value $V(S)$ for
$S \subseteq \{1,\ldots,N\}$ is the joint expected payoff when members
of~$S$ coordinate their policies.
\end{definition}

\begin{definition}[Rational coalition]
Coalition~$S$ is \emph{rational} if:
\begin{equation}\label{eq:coalition-rational}
  V(S) \;>\; \sum_{i \in S} V(\{i\}).
\end{equation}
The excess $V(S) - \sum_i V(\{i\})$ is the \emph{coordination premium}.
\end{definition}

\begin{remark}[Competitive coalitions]
Condition \eqref{eq:coalition-rational} applies in \emph{competitive}
games too. In a 6-player game, a 3-player coalition may outperform
three individuals playing against the opposing coalition. This is the
formal version of ``it pays to have allies even in competitive settings.''
\end{remark}

\subsection{Communication-Adjusted Coalition Value}

The actual coalition value depends on communication quality:
\begin{equation}
  V_{\text{comm}}(S) \;=\; V(S) - \underbrace{\sum_{i \in S} C_i \cdot L_{\text{comm}}^{(i)}}_{\text{coordination loss from imperfect communication}}
\end{equation}
where $C_i$ is the marginal coordination value of agent~$i$.

Coalition~$S$ remains rational under imperfect communication iff:
\begin{equation}
  V_{\text{comm}}(S) \;>\; \sum_{i \in S} V(\{i\}).
\end{equation}

Using the loss decomposition (Theorem~\ref{thm:loss-decomp}) and the
evidence bound (Axiom~\ref{ax:self-ev}):
\begin{equation}
  V_{\text{comm}}(S) \;\geq\; V(S) - 2\sum_{i \in S} C_i\,(V_i + L_{\text{ch}}^{(i)}).
\end{equation}

\emph{High-evidence agents (low $V_i$) are more valuable coalition members.}


\section{The Cooperative Gradient}

\subsection{Update Rule}

When agent~$i$ is in coalition~$S$, its gradient is augmented with a
cooperative correction based on communicated policies:
\begin{equation}\label{eq:coop-update}
  \pi_{i,n+1} = \proj_\Pi\Bigl(\pi_{i,n} + \gamma_n \cdot w_i \cdot
    \bigl(\hat{v}_{i,n} + \beta_n \cdot \text{coop}_i(\pi_n, \tilde{\pi}_{S \setminus i})\bigr)\Bigr)
\end{equation}
where:
\begin{itemize}
  \item $w_i = V_{\min}/V_i$ is the evidence weight (from EW-PG),
  \item $\beta_n$ is the cooperation strength schedule,
  \item $\text{coop}_i$ is the cooperative correction from communicated policies.
\end{itemize}

\begin{remark}[Triple-duty evidence weight]
The weight $w_i$ multiplies \emph{both} the standard gradient and the
cooperative term. This is correct: if agent~$i$ has poor self-knowledge
(high $V_i$, low $w_i$), it should downweight its cooperative contribution
too, since its self-model is unreliable. The evidence weight does triple duty:
\begin{enumerate}
  \item \textbf{Gradient quality}: reduces PG variance (EW-PG, Theorem~5.1 of \cite{ewpg}).
  \item \textbf{Self-knowledge}: bounds $L_{\text{self}}$ (Axiom~\ref{ax:self-ev}).
  \item \textbf{Communication quality}: bounds $L_{\text{comm}}$ (Corollary~\ref{cor:bottleneck}).
\end{enumerate}
All three are manifestations of the same underlying information content.
\end{remark}

\subsection{Communication Bias}

The cooperative term uses $\tilde{\pi}_j$ (communicated) instead of
$\pi_j$ (true). This introduces bias:
\begin{equation}
  \|\text{coop}(\pi, \pi_{S}) - \text{coop}(\pi, \tilde{\pi}_{S})\|
  \;\leq\; L_{\text{coop}} \cdot \sqrt{L_{\text{comm}}}
  \;\leq\; L_{\text{coop}} \cdot \sqrt{2V_i + 2L_{\text{ch}}}
\end{equation}
where $L_{\text{coop}}$ is the Lipschitz constant of the cooperative term.


\section{Convergence Analysis}

\subsection{Annealed Cooperation Preserves Convergence}

\begin{theorem}[Convergence of annealed cooperative PG]\label{thm:coop-conv}
Under the conditions of Giannou et al.\ (2022), with:
\begin{itemize}
  \item $\gamma_n = \gamma/(n+m)^p$, $p \in (1/2, 1]$
  \item $\beta_n = \beta/(n+m)^r$, $r > 1 - p$
  \item Cooperative term bounded: $\|\text{coop}(\pi)\| \leq K$
  \item Cooperative term vanishes at Nash: $\text{coop}(\pi^*) = 0$
\end{itemize}
the cooperative PG \eqref{eq:coop-update} converges to stable Nash
at the same rate as standard PG.
\end{theorem}

\begin{proof}[Proof sketch]
Identical structure to annealed LOLA (OpponentShapingPG, Theorem~4.3).
The cooperative term adds $\beta_n K$ to the bias bound:
$\|b_n^{\text{coop}}\| \leq B_n + \beta_n K = O(n^{-\min(\ell_b, r)})$.
The condition $r > 1 - p$ ensures $p + \min(\ell_b, r) > 1$,
so Giannou's Theorem~1 applies with the augmented bias.
\end{proof}

\subsection{Cooperative Basin Enlargement}

\begin{definition}[Cooperative reinforcement]
The cooperative Hessian $H_C = \nabla_\pi\,\text{coop}(\pi^*)$ is
\emph{cooperatively reinforcing} if its symmetric part
$S_C = (H_C + H_C^\top)/2$ is negative semi-definite.
\end{definition}

\begin{theorem}[Basin enlargement]\label{thm:coop-basin}
Under cooperative reinforcement, the cooperative PG has SOS parameter:
\begin{equation}
  \mu_{\text{coop}} = \mu + \beta \cdot \mu_C
\end{equation}
where $-\mu_C$ is the largest eigenvalue of $S_C$ (so $\mu_C \geq 0$).
The basin of attraction is strictly larger when $\mu_C > 0$.
\end{theorem}


\section{The Complete $\Omega$-Gradient}

Combining all three mechanisms, the full multi-agent gradient is:
\begin{equation}\label{eq:full-gradient}
  \boxed{
    \pi_{i,n+1} = \proj_\Pi\Bigl(\pi_{i,n} + \gamma_n \cdot w_i \cdot
    \bigl(\underbrace{\hat{v}_{i,n}}_{\text{PG}}
    + \underbrace{\lambda_n \cdot \text{OS}(\pi_n)}_{\text{LOLA}}
    + \underbrace{\beta_n \cdot \text{coop}(\pi_n, \tilde{\pi}_S)}_{\text{coalition}}\bigr)\Bigr)
  }
\end{equation}

\begin{theorem}[Full $\Omega$-PG convergence]\label{thm:full}
The update \eqref{eq:full-gradient} with annealed schedules
$\lambda_n, \beta_n \to 0$ achieves:
\begin{enumerate}
  \item \textbf{Variance improvement}: $\HM(V)/\AM(V)$ from evidence weighting.
  \item \textbf{Opponent-shaping basin}: $\mu + \lambda \mu_H$ from LOLA.
  \item \textbf{Cooperative basin}: $\mu + \beta \mu_C$ from communication.
  \item \textbf{Combined}: $\mu_{\text{full}} = \mu + \lambda \mu_H + \beta \mu_C$,
        $C_{\text{full}} = \frac{\HM}{\AM} \cdot C_{\text{std}}$.
\end{enumerate}
The three mechanisms are orthogonal: evidence weighting affects~$C$,
opponent shaping and cooperation affect~$\mu$ through independent
Hessian contributions.
\end{theorem}

\subsection{Mapping to the $\Omega$-Framework (theos eq.~27)}

The five gradient components from the $\Omega$-framework are now
fully covered:
\begin{center}
\begin{tabular}{cll}
\hline
\textbf{Component} & \textbf{$\Omega$-framework} & \textbf{Formalization} \\
\hline
(1) Exploration    & $\mathbb{E}[L \cdot s_\theta]$   & $\hat{v}_i$ (REINFORCE) \\
(2) Exploitation   & $\mathbb{E}[\nabla_\theta L]$    & $\hat{v}_i$ (backprop)  \\
(3) Evidence       & $-\mu V^{-2}\nabla V$            & $w_i$ (Keynesian weight) \\
(4) Alignment      & $\nabla D$                        & $\beta_n \cdot \text{coop}$ \\
(5) Opp.\ shaping  & $\sum_{j \neq i} \frac{\partial R_i}{\partial \theta_j} \frac{\partial \theta'_j}{\partial \theta_i}$ & $\lambda_n \cdot \text{OS}$ \\
\hline
\end{tabular}
\end{center}

Component~(4) was previously listed as ``alignment/consensus'' without
formalization. The cooperative PG fills this gap: agents align their
policies through lossy communication, with quality bounded by
evidence weights.


\section{The Six-Way Risk Decomposition --- Complete}

The $\Omega$-framework decomposes multi-agent risk into six terms:
\begin{equation}
  R_{\text{multi}} = \underbrace{R_{W \cap \Pi_N \cap B^c} + R_{S \cap \Pi_N \cap B^c}}_{\text{learnable}}
  + \underbrace{R_{W \cap \Pi_N \cap B} + R_{S \cap \Pi_N \cap B}}_{\text{G\"odel-limited}}
  + \underbrace{R_{W \cap \Pi_U} + R_{S \cap \Pi_U}}_{\text{Keynes-limited}}
\end{equation}

\begin{center}
\begin{tabular}{lll}
\hline
\textbf{Method} & \textbf{Terms addressed} & \textbf{Mechanism} \\
\hline
Standard PG     & 1--2 (learnable)     & Gradient descent \\
EW-PG           & + 5--6 (Keynes)      & Evidence-aware weighting \\
LOLA-PG         & + 3--4 (G\"odel, adv.)  & Infer opponent's blind spots \\
Coop-PG         & + 3--4 (G\"odel, coop.) & Communicate own blind spots \\
\textbf{Full $\Omega$-PG} & \textbf{All 6 terms} & \textbf{All mechanisms combined} \\
\hline
\end{tabular}
\end{center}

LOLA and cooperation address the G\"odel-limited terms through
\emph{dual} mechanisms:
\begin{itemize}
  \item \textbf{LOLA}: Agent~$i$ \emph{infers} $G_{F_j}$ from $j$'s gradient response.
        (Competitive G\"odelian step: I figure out what you can't see.)
  \item \textbf{Coop}: Agent~$j$ \emph{communicates} (a lossy version of) $G_{F_j}$.
        (Cooperative G\"odelian step: I tell you what I can see.)
\end{itemize}
These are independent information channels for the same G\"odelian content.
When both are available (intra-coalition + inter-coalition dynamics),
they provide complementary information.

Only terms in $B_{\min}$ (the irreducible collective blind spot) remain
permanently inaccessible --- as G\"odel guarantees.


\section{The ``Vibing'' Spectrum and Evolutionary Pressure}

The $\mathcal{O} \to \Pi$ mapping defines a spectrum:
\begin{equation*}
  \underbrace{\text{Full articulation}}_{\substack{L_{\text{self}} \approx 0 \\ \text{low } V_i \\ \text{Pearl Level 3}}}
  \quad\longleftrightarrow\quad
  \underbrace{\text{Pure vibing}}_{\substack{L_{\text{self}} \approx V_{\max} \\ \text{high } V_i \\ \text{Pearl Level 1}}}
\end{equation*}

The cooperative PG creates \emph{evolutionary pressure toward articulability}:
agents that can communicate their policies contribute more to coalitions,
receive higher coalition payoff, and are preferred as coalition partners.

This is the formal version of the $\Omega$-framework's claim that
multi-agent interaction drives agents up the causal hierarchy:
$\text{Pearl Level 1} \to \text{Level 2} \to \text{Level 3}$.

The pressure toward articulability is itself a \emph{phase transition}:
below a critical evidence threshold, agents cannot form effective
coalitions and are stuck in competitive dynamics. Above the threshold,
cooperation becomes possible, creating a positive feedback loop
(more cooperation $\to$ more shared information $\to$ better evidence
$\to$ better communication $\to$ more cooperation).


\section{Implications for Competitive Environments}

\subsection{Dynamic Coalition Formation}

In competitive games, coalitions are not fixed. At each stage, agents
evaluate whether coalition membership improves their payoff:
\begin{equation}
  i \in S^*_n \iff V_{\text{comm}}(S^*_n) / |S^*_n| > V(\{i\})
\end{equation}
where $S^*_n$ is the optimal coalition at step~$n$.

The communication quality (bounded by evidence) determines the
\emph{stable coalition structure}: agents with similar evidence
levels form more effective coalitions (matched communication capacity).

\subsection{The Competitive-Cooperative Duality}

The full gradient \eqref{eq:full-gradient} treats competition (LOLA)
and cooperation (coop) symmetrically:
\begin{itemize}
  \item Against opponents: LOLA term shapes their learning trajectory.
  \item With teammates: coop term aligns through communication.
  \item The evidence weight $w_i$ moderates both.
\end{itemize}

An agent simultaneously cooperates within its coalition and competes
against other coalitions. This is the generic structure of real-world
multi-agent interaction: firms cooperate internally and compete
externally; nations form alliances against common adversaries;
individuals team up in games with shifting loyalties.


\appendix
\section{Pseudocode}\label{app:algorithms}

\begin{algorithm}[H]
\caption{Cooperative PG with Lossy Communication (Coop-PG)}\label{alg:coop-pg}
\KwIn{Stochastic game $\mathcal{G}$, $N$ agents}
\KwIn{Learning rates $\gamma_n$, cooperation schedule $\beta_n = \beta/(n+m)^r$}
\KwIn{Communication channel with capacity $C$, coalition update period $T_{\text{coal}}$}
\KwOut{Policies $\{\pi_i\}_{i=1}^N$ converging to stable Nash}
\BlankLine
Initialize $\pi_i \in \Pi$ for each agent $i$\;
Initialize coalition structure $\mathcal{S} \leftarrow \{\{1\}, \{2\}, \ldots, \{N\}\}$\;
Initialize self-models $\hat{\pi}_i \leftarrow \pi_i$ for each $i$\;
\BlankLine
\For{$n = 1, 2, \ldots$}{
  \BlankLine
  \tcp{--- Phase 1: Standard gradient and evidence estimation ---}
  \For{$i = 1, \ldots, N$}{
    Collect trajectory, compute $\hat{v}_{i,n}$ and $\hat{V}_{i,n}$\;
    $w_{i,n} \leftarrow V_{\min,n} / \hat{V}_{i,n}$\;
  }
  \BlankLine
  \tcp{--- Phase 2: Self-model update ($\mathcal{O} \to \Pi$) ---}
  \For{$i = 1, \ldots, N$}{
    Update self-model from recent behavior:\;
    $\hat{\pi}_{i,n} \leftarrow \textsc{InferPolicy}(o_{i,1:n}, a_{i,1:n})$\;
    \tcp{Self-knowledge loss: $L_{\text{self}}^{(i)} \leq \hat{V}_{i,n}$}
  }
  \BlankLine
  \tcp{--- Phase 3: Communication within coalitions ---}
  \For{each coalition $S \in \mathcal{S}$}{
    \For{$i \in S$}{
      Encode self-model: $s_i \leftarrow \textsc{Encode}(\hat{\pi}_{i,n}, C)$\;
      Broadcast $s_i$ to coalition members\;
    }
    \For{$i \in S$}{
      \For{$j \in S \setminus \{i\}$}{
        Decode teammate's signal:\;
        $\tilde{\pi}_{j \to i} \leftarrow \textsc{Decode}(s_j)$\;
        \tcp{Total loss: $\leq 2 \hat{V}_{j,n} + 2 L_{\text{ch}}$}
      }
    }
  }
  \BlankLine
  \tcp{--- Phase 4: Cooperative gradient ---}
  \For{$i = 1, \ldots, N$}{
    Let $S_i$ be the coalition containing $i$\;
    \eIf{$|S_i| > 1$}{
      Compute cooperative correction from communicated policies:\;
      $\text{coop}_{i,n} \leftarrow \nabla_{\theta_i} \hat{R}_{S_i}(\pi_{i,n},\, \tilde{\pi}_{S_i \setminus i})$\;
    }{
      $\text{coop}_{i,n} \leftarrow 0$\;
    }
    \BlankLine
    $\pi_{i,n+1} \leftarrow \proj_\Pi\!\bigl(\pi_{i,n} + \gamma_n \cdot w_{i,n} \cdot (\hat{v}_{i,n} + \beta_n \cdot \text{coop}_{i,n})\bigr)$\;
  }
  \BlankLine
  \tcp{--- Phase 5: Coalition update (every $T_{\text{coal}}$ steps) ---}
  \If{$n \bmod T_{\text{coal}} = 0$}{
    \For{each pair $i, j$ in different coalitions}{
      Estimate coalition value $\hat{V}(S_i \cup S_j)$ from recent data\;
      \If{$\hat{V}(S_i \cup S_j) > \hat{V}(S_i) + \hat{V}(S_j)$}{
        Merge: $\mathcal{S} \leftarrow \mathcal{S} \setminus \{S_i, S_j\} \cup \{S_i \cup S_j\}$\;
      }
    }
    \For{each coalition $S$ with $|S| > 1$}{
      \For{$i \in S$}{
        \If{$V(\{i\}) > V(S) / |S|$}{
          Split: agent $i$ leaves $S$\;
        }
      }
    }
  }
}
\end{algorithm}

\begin{algorithm}[H]
\caption{Full $\Omega$-PG (EW-LOLA-Coop)}\label{alg:full-omega}
\KwIn{All inputs from EW-PG, LOLA-PG, Coop-PG}
\KwIn{Shaping schedule $\lambda_n$, cooperation schedule $\beta_n$}
\KwOut{Policies converging to stable Nash with all three improvements}
\BlankLine
Initialize policies, variances, opponent models, self-models, coalitions\;
\BlankLine
\For{$n = 1, 2, \ldots$}{
  \BlankLine
  \tcp{Compute all components}
  $\hat{v}_{i,n},\; w_{i,n}$ for each $i$ \tcp*{EW-PG (gradient + evidence)}
  $\text{OS}_{i,n}$ for each $i$ \tcp*{LOLA (opponent shaping)}
  $\tilde{\pi}_{j \to i}$ for $j \in S_i$ \tcp*{Communication}
  $\text{coop}_{i,n}$ for each $i$ \tcp*{Cooperative correction}
  \BlankLine
  \tcp{Full $\Omega$-update (theos eq.~27, all 5 components)}
  \For{$i = 1, \ldots, N$}{
    $g_{i,n} \leftarrow \hat{v}_{i,n} + \lambda_n \cdot \text{OS}_{i,n} + \beta_n \cdot \text{coop}_{i,n}$\;
    $\pi_{i,n+1} \leftarrow \proj_\Pi\!\bigl(\pi_{i,n} + \gamma_n \cdot w_{i,n} \cdot g_{i,n}\bigr)$\;
  }
  \BlankLine
  Update coalitions every $T_{\text{coal}}$ steps \tcp*{Phase 5 of Alg.~\ref{alg:coop-pg}}
}
\BlankLine
\tcp{Convergence guarantees (Theorem~\ref{thm:full}):}
\tcp{$\quad \E[\|\pi_n - \pi^*\|^2] = O\bigl(\frac{\HM(V)}{\AM(V)} \cdot C_{\text{std}}\, /\, n^q\bigr)$}
\tcp{$\quad \mu_{\text{full}} = \mu + \lambda \mu_H + \beta \mu_C$}
\end{algorithm}

\begin{remark}[Implementation notes]
\begin{enumerate}
  \item \textsc{InferPolicy} reconstructs $\hat{\pi}_i$ from observation-action
    history. Options: (a)~frequency counts (tabular), (b)~behavioral cloning
    network trained on own $(s,a)$ pairs, (c)~inverse RL.
  \item \textsc{Encode}/\textsc{Decode} implement lossy compression of the
    self-model. For continuous policies: quantize and transmit parameters.
    For tabular: transmit action distributions with precision $\propto C$.
  \item Coalition updates (Phase~5) are expensive. The period $T_{\text{coal}}$
    should be large enough for value estimates to stabilize (e.g., 100--1000 steps).
  \item In competitive games, coalitions are temporary. An agent defects when
    its individual value exceeds its per-capita coalition share.
    The split check (Phase~5) handles this.
  \item The self-knowledge loss $L_{\text{self}}^{(i)} \leq \hat{V}_{i,n}$
    is monitored but not explicitly minimized. It decreases naturally as
    the agent gathers more evidence. To accelerate: add an intrinsic reward
    for self-model accuracy (curiosity about own policy).
\end{enumerate}
\end{remark}

\bibliographystyle{plain}
\bibliography{references}

\end{document}
